\subsection{Rank dan Sistem Homogen}

Terdapat jenis sistem khusus yang memerlukan kajian tambahan. Jenis sistem ini disebut sistem persamaan
homogen, yang telah kita definisikan dalam Definisi \ref{def:homogeneoussystem} di atas.
Fokus kita dalam bagian ini adalah membahas jenis-jenis solusi yang mungkin bagi suatu sistem persamaan homogen.

Perhatikan definisi berikut.

\begin{definition}{Solusi Trivial}\label{def:trivialsolution}
Perhatikan sistem persamaan homogen yang diberikan oleh
\begin{equation*}
\begin{array}{c}
a_{11}x_{1}+a_{12}x_{2}+\cdots +a_{1n}x_{n}= 0 \\
a_{21}x_{1}+a_{22}x_{2}+\cdots +a_{2n}x_{n}= 0  \\
\vdots \\
a_{m1}x_{1}+a_{m2}x_{2}+\cdots +a_{mn}x_{n}= 0 
\end{array}
\end{equation*}
Maka $x_{1} = 0, x_{2} = 0, \cdots, x_{n} =0$ selalu merupakan
solusi sistem ini. Kita menyebutnya \textbf{solusi trivial} \index{solusi trivial}.
\end{definition}

Jika sistem memiliki solusi yang tidak membuat semua $x_1, \cdots, x_n$ sama dengan nol,
kita menyebut solusi tersebut \textbf{nontrivial} \index{solusi nontrivial}. Solusi trivial
tidak memberi banyak informasi tentang sistem karena hanya menyatakan bahwa $0=0$.
Oleh karena itu, ketika bekerja dengan sistem persamaan
homogen, kita ingin mengetahui kapan sistem tersebut memiliki solusi nontrivial.

Misalkan kita memiliki suatu sistem homogen yang terdiri atas $m$ persamaan dan menggunakan $n$ variabel, serta misalkan
$n > m$. Dengan kata lain, terdapat lebih banyak variabel daripada persamaan.
Ternyata sistem ini selalu memiliki solusi nontrivial. Sistem tersebut bukan saja
memiliki solusi nontrivial, melainkan juga memiliki tak berhingga banyak solusi.
Solusi nontrivial juga mungkin, tetapi tidak harus, ada jika $n=m$ dan $n<m$.

Perhatikan contoh berikut.

\begin{example}{Solusi Sistem Persamaan Homogen}\label{exa:homogeneoussolution}
Temukan solusi nontrivial sistem persamaan homogen berikut
\begin{equation*}
\begin{array}{cccccccc}
2x &+& y &-& z &=& 0 \\
x &+& 2y &-& 2z &=& 0
\end{array}
\end{equation*}
\end{example}

\begin{solution}
Perhatikan bahwa sistem ini memiliki $m = 2$ persamaan dan $n = 3$ variabel, sehingga $n>m$.
Oleh karena itu, berdasarkan pembahasan sebelumnya, kita memperkirakan bahwa sistem ini memiliki tak berhingga banyak solusi.

Proses yang kita gunakan untuk menemukan solusi sistem persamaan homogen
sama dengan proses yang digunakan pada bagian sebelumnya.
Pertama, kita menyusun matriks diperbesar, yaitu
\begin{equation*}
\leftB
\begin{array}{rrr|r}
2 & 1 & -1 & 0 \\ 
1 & 2 & -2 & 0
\end{array}
\rightB
\end{equation*}
Kemudian, kita membawa matriks ini ke \rref-nya, yang diberikan di bawah.
\begin{equation*}
\leftB
\begin{array}{rrr|r}
1 & 0 & 0 & 0 \\ 
0 & 1 & -1 & 0
\end{array}
\rightB
\end{equation*}
Sistem persamaan yang bersesuaian adalah
\begin{equation*}
\begin{array}{c}
x = 0 \\
y - z =0 \\
\end{array}
\end{equation*}
Karena $z$ tidak dibatasi oleh persamaan mana pun, kita mengetahui bahwa variabel ini akan menjadi parameter.
Misalkan $z=t$, dengan $t$ berupa sembarang bilangan.
Oleh karena itu, solusi kita berbentuk
\begin{equation*}
\begin{array}{c}
x = 0 \\
y = z = t \\
z = t
\end{array}
\end{equation*}
Jadi, sistem ini memiliki tak berhingga banyak solusi dengan satu parameter $t$.
\end{solution}

Misalkan kita menuliskan solusi contoh sebelumnya dalam bentuk lain.
Secara khusus,
\begin{equation*}
\begin{array}{c}
x = 0 \\
y = 0 + t \\
z = 0 + t
\end{array}
\end{equation*}
dapat ditulis sebagai
\begin{equation*}
\leftB
\begin{array}{r}
x\\
y\\
z
\end{array}
\rightB
 =
\leftB
\begin{array}{r}
0\\
0\\
0
\end{array}
\rightB
 +
t
\leftB
\begin{array}{r}
0\\
1\\
1
\end{array}
\rightB
\end{equation*}
Perhatikan bahwa kita telah menyusun sebuah kolom dari konstanta-konstanta dalam solusi (semuanya sama dengan $0$),
 serta sebuah kolom yang bersesuaian dengan
koefisien $t$ dalam setiap persamaan. Meskipun bentuk solusi ini akan dibahas lebih lanjut pada bab-bab berikutnya,
untuk saat ini perhatikan kolom koefisien
parameter $t$. Dalam hal ini, kolom tersebut adalah
$\leftB
\begin{array}{r}
0\\
1\\
1
\end{array}
\rightB$. 

Kolom ini memiliki nama khusus, yaitu \textbf{solusi dasar}\index{solusi dasar}. Solusi dasar suatu sistem
adalah kolom yang disusun dari koefisien parameter dalam solusi. Kita sering menyatakan solusi dasar
dengan $X_1, X_2$, dan seterusnya, bergantung pada banyaknya solusi yang muncul. Oleh karena itu, Contoh \ref{exa:homogeneoussolution}
memiliki solusi dasar $X_1 = \leftB
\begin{array}{r}
0\\
1\\
1
\end{array}
\rightB$. 

Kita membahasnya lebih lanjut dalam contoh berikut.

\begin{example}{Solusi Dasar Sistem Homogen}\label{exa:basicsolutions}
Perhatikan sistem persamaan homogen berikut.
\begin{equation*}
\begin{array}{cccccccc}
x &+& 4y &+& 3z &=& 0 \\
3x &+& 12y &+& 9z &=& 0
\end{array}
\end{equation*}
Temukan solusi dasar sistem ini.
\end{example}

\begin{solution}
Matriks diperbesar sistem ini dan \rref \;yang dihasilkan adalah
\begin{equation*}
\leftB
\begin{array}{rrr|r}
1 & 4 & 3 & 0 \\
3 & 12 & 9 & 0
\end{array}
\rightB
\rightarrow \cdots \rightarrow
\leftB
\begin{array}{rrr|r}
1 & 4 & 3 & 0 \\
0 & 0 & 0 & 0
\end{array}
\rightB
\end{equation*}
Jika ditulis sebagai persamaan, sistem ini diberikan oleh
\begin{equation*}
x + 4y +3z=0
\end{equation*}
Perhatikan bahwa hanya $x$ yang bersesuaian dengan kolom pivot. Dalam hal ini, kita akan memiliki dua parameter,
satu untuk $y$ dan satu untuk $z$. Misalkan $y = s$ dan $z=t$ untuk sembarang bilangan $s$ dan $t$. Maka solusi kita menjadi
\begin{equation*}
\begin{array}{c}
x = -4s - 3t \\
y = s \\
z = t
\end{array}
\end{equation*}
yang dapat ditulis sebagai
\begin{equation*}
\leftB
\begin{array}{r}
x\\
y\\
z
\end{array}
\rightB
=
\leftB
\begin{array}{r}
0\\
0\\
0
\end{array}
\rightB
+
s
\leftB
\begin{array}{r}
-4 \\
1 \\
0
\end{array}
\rightB
+ 
t
\leftB
\begin{array}{r}
-3 \\
0 \\
1
\end{array}
\rightB
\end{equation*}
Di sini Anda dapat melihat bahwa kita memiliki dua kolom koefisien yang bersesuaian dengan parameter, yaitu satu untuk $s$ dan satu untuk $t$.
Oleh karena itu, sistem ini memiliki dua solusi dasar, yaitu
\begin{equation*}
X_1=
\leftB
\begin{array}{r}
-4 \\
1 \\
0
\end{array}
\rightB, X_2 = \leftB
\begin{array}{r}
-3 \\
0 \\
1
\end{array}
\rightB
\end{equation*} 
\end{solution}

Sekarang kita menyajikan definisi baru.

\begin{definition}{Kombinasi Linear}\label{def:linearcombination-columns}
Misalkan $X_1,\cdots ,X_n,V$ merupakan matriks kolom. Maka
$V$ disebut \textbf{kombinasi linear}
\index{kombinasi linear}dari kolom-kolom $X_1,\cdots , X_n $
jika terdapat skalar $a_{1},\cdots ,a_{n}$ sedemikian
sehingga
\begin{equation*}
V = a_1 X_1 + \cdots + a_n X_n
\end{equation*}
\end{definition}

Salah satu hasil penting bagian ini adalah bahwa kombinasi linear dari solusi-solusi dasar kembali merupakan solusi sistem.
Yang lebih menarik lagi, setiap solusi dapat ditulis sebagai kombinasi linear dari solusi-solusi tersebut.
Oleh karena itu, jika kita mengambil kombinasi linear dari kedua solusi pada Contoh \ref{exa:basicsolutions},
hasilnya juga merupakan solusi.
Sebagai contoh, kita dapat mengambil kombinasi linear berikut
\begin{equation*}
3
\leftB
\begin{array}{r}
-4 \\
1 \\
0
\end{array}
\rightB
+
2
\leftB
\begin{array}{r}
-3 \\
0\\
1
\end{array}
\rightB
 =
\leftB
\begin{array}{r}
-18 \\
3 \\
2
\end{array}
\rightB
\end{equation*}
Luangkan waktu sejenak untuk memverifikasi bahwa
\begin{equation*}
\leftB
\begin{array}{r}
x \\
y \\
z
\end{array}
\rightB
=
\leftB
\begin{array}{r}
-18 \\
3 \\
2
\end{array}
\rightB
\end{equation*}
memang merupakan solusi sistem pada Contoh \ref{exa:basicsolutions}.

Cara lain untuk memperoleh lebih banyak informasi mengenai solusi
suatu sistem homogen adalah dengan memperhatikan \textbf{rank} matriks koefisien yang bersesuaian. Sekarang kita
mendefinisikan arti rank suatu matriks.

\begin{definition}{Rank Suatu Matriks}\label{def:rank}
Misalkan $A$ merupakan suatu matriks dan perhatikan sembarang \ef \;dari $A$.
Maka banyaknya $r$ entri utama $A$ tidak bergantung pada \ef \;yang Anda pilih dan disebut \textbf{rank} \index{matriks!rank} dari $A$. Kita menyatakannya dengan $\func{rank}(A)$.
\end{definition}

Dengan cara serupa, kita dapat menghitung banyaknya posisi pivot (atau kolom pivot)
untuk menentukan rank $A$.

\begin{example}{Menentukan Rank Suatu Matriks}\label{exa:rankofamatrix}
Perhatikan matriks
\begin{equation*}
\leftB
\begin{array}{rrr}
1 & 2 & 3 \\
1 & 5 & 9 \\
2 & 4 & 6 
\end{array}
\rightB
\end{equation*}
Berapakah rank-nya?
\end{example}

\begin{solution}
Pertama, kita perlu menemukan \rref \;dari $A$. Dengan algoritme yang biasa, kita memperoleh
\begin{equation*}
\leftB
\begin{array}{rrr}
\fbox{1} & 0 & -1 \\
0 & \fbox{1} & 2 \\
0 & 0 & 0
\end{array}
\rightB
\end{equation*}
Di sini kita memiliki dua entri utama, atau dua posisi pivot, yang diperlihatkan dalam kotak di atas. Rank $A$ adalah $r = 2.$
\end{solution}

Perhatikan bahwa kita akan memperoleh jawaban yang sama jika kita menemukan \ef \;dari $A$, bukan \rref-nya.

Misalkan kita memiliki suatu sistem homogen yang terdiri atas $m$ persamaan dalam $n$ variabel, dan misalkan
$n > m$. Dari pembahasan di atas, kita mengetahui bahwa sistem ini memiliki tak berhingga banyak solusi. Jika kita memperhatikan rank matriks
 koefisien sistem ini, kita dapat memperoleh lebih banyak informasi mengenai solusinya. Perhatikan bahwa kita hanya melihat matriks koefisien, bukan seluruh
matriks diperbesar.

\begin{theorem}{Rank dan Solusi Sistem Homogen}\label{thm:rankhomogeneoussolutions}
Misalkan $A$ merupakan matriks koefisien $m \times n$ yang bersesuaian dengan suatu sistem persamaan homogen, dan misalkan $A$ memiliki rank $r$.
Maka solusi sistem yang bersesuaian memiliki $n-r$ parameter.
\end{theorem}

Perhatikan Contoh \ref{exa:basicsolutions} di atas dalam konteks
teorema ini. Sistem dalam contoh tersebut memiliki $m = 2$ persamaan dalam $n =
3$ variabel. Pertama, karena $n>m$, kita mengetahui bahwa sistem tersebut memiliki
solusi nontrivial, dan karena itu memiliki tak berhingga banyak solusi. Hal ini
memberi tahu kita bahwa solusi akan memuat sedikitnya satu parameter. Rank
matriks koefisien dapat memberikan lebih banyak informasi mengenai
solusi. Rank matriks koefisien sistem tersebut adalah $1$ karena
matriks itu memiliki satu entri utama dalam \ef. Teorema \ref{thm:rankhomogeneoussolutions} memberi tahu kita bahwa
solusi akan memiliki $n-r = 3-1 = 2$ parameter. Anda dapat memeriksa bahwa hal ini
benar dalam solusi Contoh \ref{exa:basicsolutions}.

Perhatikan bahwa jika $n=m$ atau $n<m$, sistem dapat memiliki solusi unik (yang merupakan solusi trivial) atau tak berhingga banyak solusi.

Di sini kita tidak terbatas pada sistem persamaan homogen. Rank suatu matriks dapat digunakan untuk memperoleh informasi mengenai
solusi sembarang sistem persamaan linear.
Pada bagian sebelumnya, kita membahas bahwa suatu sistem persamaan dapat tidak memiliki solusi, memiliki solusi unik, atau memiliki tak berhingga banyak solusi.
Misalkan sistem tersebut konsisten, baik homogen maupun tidak. Teorema berikut memberi tahu kita
cara menggunakan rank untuk mengetahui jenis solusi yang kita miliki.

\begin{theorem}{Rank dan Solusi Sistem Persamaan Konsisten}\label{thm:rankconsistentsolutions}
Misalkan $A$ merupakan matriks diperbesar $m \times \left( n+1 \right) $ yang bersesuaian dengan suatu sistem persamaan konsisten dalam $n$ variabel, dan misalkan $A$ memiliki rank $r$.
Maka
\begin{enumerate}
\item sistem memiliki solusi unik jika $r = n $
\item sistem memiliki tak berhingga banyak solusi jika $r < n$
\end{enumerate}
\end{theorem}

Kita tidak akan menyajikan pembuktian formal untuk hasil ini, tetapi perhatikan pembahasan berikut.

\begin{enumerate}
\item {\em Tidak Ada Solusi \em}
Teorema di atas mengasumsikan bahwa sistem tersebut konsisten, yaitu bahwa sistem memiliki solusi. Ternyata matriks
diperbesar dari suatu sistem tanpa solusi dapat memiliki sembarang rank $r$ asalkan $r>1$. Oleh karena itu, kita harus mengetahui
bahwa sistem tersebut konsisten agar dapat menggunakan teorema ini.

\item {\em Solusi Unik \em}
Misalkan $r=n$. Maka terdapat posisi pivot pada setiap kolom matriks koefisien $A$. Jadi, terdapat solusi unik.

\item {\em Tak Berhingga Banyak Solusi \em}
Misalkan $r<n$. Maka terdapat tak berhingga banyak solusi. Banyaknya posisi pivot (dan karena itu banyaknya entri utama) lebih kecil daripada banyaknya kolom, yang berarti bahwa tidak setiap kolom merupakan kolom pivot. Kolom-kolom yang tidak merupakan kolom pivot bersesuaian dengan parameter. Bahkan,
dalam hal ini kita memiliki $n-r$ parameter.
\end{enumerate}
